\bookchapter{Chapitre 54 : Projet 1 — Mon carnet de comptes (version locale)}

\chapterintro{Bienvenue dans ton premier « vrai projet ». Ce carnet de comptes ne se connecte pas au réseau, mais il a tout d'une grande : enregistrer une écriture, consulter l'historique, faire des statistiques mensuelles, supprimer et modifier. Une fois terminé, tu auras parcouru de bout en bout le flux complet « données logique interface ». Ne te précipite pas pour recopier le code — réfléchis d'abord, agis ensuite : c'est notre principe de toujours.}

\booksection{Réfléchir au besoin : pose d'abord cinq questions}

Avant de te lancer, clarifie ton idée avec les « cinq éléments du besoin » :

\begin{center}
\begin{tabularx}{\linewidth}{|X|X|}
\hline
\textbf{Élément} & \textbf{Réponse} \tabularnewline
\hline
Who (pour qui) & Pour soi, et pour la famille ou les amis aussi \tabularnewline
What (faire quoi) & Enregistrer revenus et dépenses, consulter l'historique, résumer par mois, statistiques par catégorie \tabularnewline
Why (pourquoi) & S'entraîner + comprendre où part l'argent \tabularnewline
How (comment) & Kotlin + Compose + Room + MVVM, 100 \% local, sans back-end \tabularnewline
Constraint (contraintes) & Terminer en une semaine ; Android 8 minimum ; même sans publication, coder selon les critères de publication \tabularnewline
\hline
\end{tabularx}
\end{center}

\begin{center}\small\color{gray}Les cinq éléments du besoin du carnet de comptes\end{center}

Décompose ensuite en petites étapes « utilisables dès leur achèvement » : (1) la couche données Room (stocker une écriture) ; (2) l'écran liste (voir l'historique) ; (3) l'écran de saisie (ajouter une écriture) ; (4) suppression et édition ; (5) la carte de synthèse mensuelle. À chaque étape finie, l'app est déjà « quelque chose d'utilisable ».

\booksection{La couche données : donner un « foyer » aux écritures}

Définis d'abord à quoi ressemble une écriture — des champs suffisants, sans gourmandise :

\begin{lstlisting}[language=Kotlin]
@Entity(tableName = "records")
data class Record(
    @PrimaryKey(autoGenerate = true) val id: Long = 0,
    val type: String,        // "income" revenu | "expense" dépense
    val amount: Double,      // montant (yuan)
    val category: String,    // catégorie : repas, transport, salaire…
    val note: String = "",   // note (facultatif)
    val date: String         // "2026-08-15"
)
\end{lstlisting}

Ajoute dans le Dao deux requêtes « pour les statistiques » — c'est l'âme du carnet de comptes :

\begin{lstlisting}[language=Kotlin]
@Dao
interface RecordDao {
    @Query("SELECT * FROM records ORDER BY date DESC, id DESC")
    fun getAll(): Flow<List<Record>>

    @Insert
    suspend fun insert(record: Record)

    @Delete
    suspend fun delete(record: Record)

    @Query("SELECT COALESCE(SUM(amount),0) FROM records WHERE type='income' AND date LIKE :month || '%'")
    suspend fun monthIncome(month: String): Double   // month est au format "2026-08"

    @Query("SELECT COALESCE(SUM(amount),0) FROM records WHERE type='expense' AND date LIKE :month || '%'")
    suspend fun monthExpense(month: String): Double
}
\end{lstlisting}

\begin{tcbinfo}{Petit conseil : pourquoi Double pour les montants ?}
Les apps financières rigoureuses utilisent BigDecimal ou un « Int en centimes » pour éviter les erreurs de virgule flottante (le vieux problème 0,1+0,2$\neq$0,3). Pour un projet d'entraînement, Double suffit ; le jour où tu feras un vrai produit financier, pense à passer aux centimes comme unité — c'est un piège que d'innombrables aînés ont foulé avant toi.

\end{tcbinfo}

\booksection{MVVM : quand interface et logique « quittent la maison commune »}

Laisser l'interface toucher directement à la base de données ? Ça tourne, mais le code de l'interface devient long et pénible. L'architecture MVVM donne à chacun son rôle : la View (l'interface) ne fait que dessiner, le ViewModel (le pivot logique) gère les données, le Model (Room) gère le stockage.

Vue (interface Compose)

ViewModel (état et logique)

Model (couche de données Room)

\begin{center}\small\color{gray}Figure 5 : la structure MVVM en trois couches — les données circulent dans un seul sens, chaque couche reste à sa place\end{center}

\begin{lstlisting}[language=Kotlin]
class RecordViewModel(private val dao: RecordDao) : ViewModel() {

    // convertit le Flow en état observable par l'interface
    val records = dao.getAll()
        .stateIn(viewModelScope, SharingStarted.WhileSubscribed(), emptyList())

    // actions utilisateur : appeler des fonctions suspend
    fun addRecord(type: String, amount: Double, category: String, note: String, date: String) {
        viewModelScope.launch {
            dao.insert(Record(type = type, amount = amount,
                category = category, note = note, date = date))
        }
    }

    fun deleteRecord(record: Record) {
        viewModelScope.launch { dao.delete(record) }
    }
}
\end{lstlisting}

Dans l'interface (Compose), écris : val vm: RecordViewModel = viewModel(), puis val list by vm.records.collectAsState() — dès que les données changent, la liste se rafraîchit toute seule, sans que tu aies rien à faire. C'est la puissance des trois mousquetaires MVVM + Flow + Compose réunis.

\booksection{L'interface : l'écran liste + l'écran de saisie}

Pour l'écran liste, utilise LazyColumn (liste à chargement paresseux) — jamais de ralentissement, même avec beaucoup de données :

\begin{lstlisting}[language=Kotlin]
@Composable
fun RecordListScreen(vm: RecordViewModel) {
    val records by vm.records.collectAsState()

    Column {
        // carte de synthèse en haut (revenus - dépenses = solde)
        SummaryCard(records)

        // liste des opérations
        LazyColumn {
            items(records, key = { it.id }) { record ->
                RecordRow(record) { vm.deleteRecord(record) }
            }
        }

        // bouton flottant "ajouter une dépense"  ouvre le formulaire
        FloatingActionButton(onClick = { showAddDialog = true }) {
            Text("+")
        }
    }
}

@Composable
fun RecordRow(record: Record, onDelete: () -> Unit) {
    Row(Modifier.fillMaxWidth().padding(12.dp)) {
        Text(record.category, fontWeight = FontWeight.Bold)
        Spacer(Modifier.weight(1f))
        Text(if (record.type == "income") "+${record.amount}" else "-${record.amount}",
             color = if (record.type == "income") Color(0xFF27AE60) else Color(0xFFC0392B))
    }
}
\end{lstlisting}

Pour la saisie, un AlertDialog (boîte de dialogue) est le plus simple : quelques champs + un basculement dépense/revenu + un menu déroulant de catégories, et au clic sur OK on appelle vm.addRecord(...).

\begin{tcbwarn}{Attention : ne fais pas l'économe sur la validation des saisies}
Le champ du montant doit absolument être validé : valeurs vides, négatives, saisies illégales comme « abc », tout doit être intercepté (avec toDoubleOrNull()). L'une des causes les plus courantes de plantage d'une app, c'est de « considérer la saisie de l'utilisateur comme un enfant sage ». Retiens : tout ce que tape l'utilisateur doit être inspecté comme un ennemi.

\end{tcbwarn}

\booksection{Éditer et supprimer : permettre aux données « d'entrer et de sortir »}

Un carnet qui sait ajouter mais pas supprimer est bon à jeter au bout de trois jours. Ajoute deux opérations :

\begin{lstlisting}[language=Kotlin]
// ajoute deux méthodes au Dao :
@Query("UPDATE records SET amount=:amount, category=:category, note=:note WHERE id=:id")
suspend fun update(id: Long, amount: Double, category: String, note: String)

@Query("DELETE FROM records WHERE id=:id")
suspend fun deleteById(id: Long)

// à l'écran : appui long sur un élément ouvre le menu "modifier/supprimer" (avec Compose : DropdownMenu ou AlertDialog)
// avant de supprimer, mieux vaut une confirmation : la "suppression accidentelle" est ce que les utilisateurs d'apps de dépenses détestent le plus
\end{lstlisting}

Deux leçons de détail sur l'interaction : la suppression exige une double confirmation (contre les doigts qui glissent), et l'édition pré-remplit autant que possible les valeurs d'origine (corriger un mot est plus simple que tout retaper). Ne sous-estime pas ces deux détails — le « côté produit » se bâtit sur ce genre de détails.

\booksection{Réception : trois niveaux de vérification + deux défis avancés}

\booksection{Réception : trois niveaux de vérification + deux défis avancés}

Projet terminé, fais la réception selon trois niveaux « fonctionnalités cas limites architecture » : sait-on enregistrer, sait-on supprimer, les données sont-elles encore là à la réouverture de l'app (fonctionnalités) ; l'app plante-t-elle avec un montant à 0 ou négatif, les données se perdent-elles à la rotation de l'écran (cas limites) ; interface et données sont-elles mélangées, des champs sont-ils codés en dur (architecture).

Une fois la version de base réceptionnée, offre-toi deux « plats avancés » :

\begin{enumerate}
  \item Filtre mensuel : ajoute en haut un sélecteur de mois (par exemple 2026-08 v), la liste n'affichant que les écritures du mois ;
  \item Camembert par catégorie : avec le Canvas de Compose ou une petite bibliothèque de graphiques, dessine « où part l'argent » en camembert — une fois fait, le carnet de comptes prend un petit « air de produit ».
\end{enumerate}

\begin{tcbpractice}{À toi de jouer : termine-le cette semaine}
Avance dans l'ordre « couche données ViewModel écran liste écran de saisie carte de synthèse », une étape par jour. En cas de blocage, balance le message d'erreur tel quel à un moteur de recherche ou à une IA (tu te souviens de la partie 7, « L'IA au service du code » ? Décris le problème en précisant « comportement attendu, comportement observé, étapes de reproduction »). Une fois l'app finie, note tes comptes pendant une semaine — utiliser une app de ta propre fabrication donne une fierté qui te poussera à continuer.

\end{tcbpractice}

\begin{tcbsummary}{Résumé du chapitre}
\begin{itemize}
  \item Cinq éléments du besoin décomposition en petites étapes « utilisables à chaque pas », puis on se lance
  \item Couche données : entité Record + Dao (avec le SQL de statistiques mensuelles) + Database
  \item Répartition MVVM : la View ne fait que dessiner, le ViewModel gère la logique, le Model gère le stockage
  \item Flow + stateIn + collectAsState : dès que les données changent, l'interface se rafraîchit automatiquement
  \item LazyColumn pour la liste à chargement paresseux ; AlertDialog pour la saisie
  \item Réception en trois niveaux (fonctionnalités/cas limites/architecture) + défis avancés (filtre mensuel, camembert par catégorie)
\end{itemize}

\end{tcbsummary}
